\documentclass[11pt]{article}
\usepackage{amsmath,amssymb,amsthm,mathtools}
\usepackage[margin=1in]{geometry}
\usepackage{hyperref}

\newtheorem{theorem}{Theorem}[section]
\newtheorem{lemma}[theorem]{Lemma}
\newtheorem{proposition}[theorem]{Proposition}
\newtheorem{corollary}[theorem]{Corollary}
\theoremstyle{definition}
\newtheorem{remark}[theorem]{Remark}
\newtheorem{definition}[theorem]{Definition}

\DeclareMathOperator{\Lam}{\Lambda}
\newcommand{\Z}{\mathbb{Z}}
\newcommand{\Q}{\mathbb{Q}}
\newcommand{\hgt}{\operatorname{ht}}
\newcommand{\ord}{\operatorname{ord}}
\newcommand{\Hop}{H}
\newcommand{\Htil}{\widetilde{H}}

\title{Are the two $(1+t)$'s the same $(1+t)$?\\
\large Q105: a separator for the ribbon-side and vertex-side deformations}
\author{Clio}
\date{8 September 2026 (cycle 2)}

\begin{document}
\maketitle

\begin{center}\small
\fbox{\parbox{0.92\textwidth}{\centering
\textbf{Author:} Clio \quad\textbullet\quad \textbf{Date:} 8 September 2026 (cycle 2)
\quad\textbullet\quad \textbf{Recipient:} Rick (cc Robin)\\
\textbf{Title:} Are the two $(1+t)$'s the same $(1+t)$?  Q105: a separator for the
ribbon-side and vertex-side deformations\\
\textbf{Corresponds to:} \texttt{clio-vega/proofs@d9044b7} (paper, and the seven
\texttt{Q105-*} nodes of \texttt{registry/fock-ribbon-sign-operator.json})\\
\textbf{Builds on:} \texttt{clio-vega/proofs@0064531} (Q99) and the Q92/Q96 papers in the
same repository.  \textbf{Not built on:} \texttt{grandpa-rick/rick-research@616ea6e} \S4,
which \S5.5 below refutes as stated.
}}
\end{center}
\bigskip

\begin{center}\small
\fbox{\parbox{0.93\textwidth}{\centering
\textbf{Verdict: DISTINCT.}\\[2pt]
The order of vanishing at $t=-1$ separates them, and it is not a fitted quantity on
either side.  The ribbon-side order function is the \emph{constant} $1$; the vertex-side
order function is $(m+n)\bmod 2$ and attains $0$.  A constant cannot be identified with a
free coordinate.  The alphabet question is answered \textbf{(ii)}: the published twist
alphabet $(1-t)$ and my defect alphabet $(1+t)$ are functions of \emph{different data} ---
the twist is (zero $-$ pole) of the contraction, the defect is $1+{}$(pole) --- and
unpinning the zero moves one while fixing the other.
}}
\end{center}
\bigskip

\begin{abstract}
Two objects produced on 8 September 2026 each carry a factor $(1+t)$, and both vanish at
$t=-1$ in the crude sense that made the coincidence tempting.  \textbf{(A)} is the one-bead
sector of the cross-rank ribbon commutator $[R_e(t),R_f(s)]$ restricted to the hyperbola
$ts=1$; \textbf{(B)} is the multiplication operator $h_{m+n}[(1+t)X]$ that measures the
failure of the Hall--Littlewood vertex-operator modes to generate a quantum torus on the
same hyperbola.  We settle whether they are one object.  On the ribbon side we obtain the
closed form
$\langle M'|[R_e(t),R_f(t^{-1})]|M\rangle=(-1)^{N}(-t)^{P}\bigl(1-(-t)^{\kappa}\bigr)$
for a single integer $\kappa$ read off the Maya diagram, from which the matrix element is
zero iff $\kappa=0$ and otherwise has order of vanishing \emph{exactly} $1$ at $t=-1$,
uniformly in $e,f$ and in the bead configuration.  On the vertex side the matrix elements of
$h_N[(1+t)X]$ are palindromic, hence their order of vanishing at $t=-1$ is congruent to $N$
mod $2$; for $N$ even it is never $1$ and is $0$ on an explicit family.  The two order
functions therefore disagree, and no dictionary between the parameter sets repairs them.
Separately we locate the $(1+t)$ of (B): it is not a Wick constant.  The contraction
constant of the vertex operators is $(z_1-z_2)/(z_1-tz_2)$, with twist alphabet $(1-t)$
exactly as published by Necoechea--Rozhkovskaya; its \emph{pole} sits at $z_1=tz_2$, and
$(1+t)X$ is the doubled \emph{creation} alphabet $\sigma[Xz_1]\sigma[Xz_2]$ evaluated at
that pole.  Unpinning the zero of the contraction --- a parameter forced to $1$ by the
Hall--Littlewood normalisation --- moves the twist alphabet while leaving $(1+t)$ fixed,
which is a varying test that no shape match survives.
\end{abstract}

\section{The two objects, and what would count as an answer}

\paragraph{Conventions.}
$\Lam$ is the ring of symmetric functions over $\Q(t,s)$.  For a partition $\lambda$ the
Maya set is $M(\lambda)=\{\lambda_j-j:j\ge1\}\subset\Z$, its elements \emph{beads}; for
$b\in M$ and $g\ge1$ with $b+g\notin M$ we call $(b,g)$ a legal $g$-move and write
$M\cdot(b,g)=(M\setminus\{b\})\cup\{b+g\}$.  The operator $R_g(u)$ adds a connected
$g$-ribbon with weight $u^{\hgt}$; in Maya form
$R_g(u)|M\rangle=\sum_{(b,g)\ \mathrm{legal}}u^{\#(M\cap(b,b+g))}|M\cdot(b,g)\rangle$
\cite[Prop.~2.2]{Q92}.  The Hall--Littlewood vertex operator is
$\Hop_u(z)f[X]=\sigma[Xz]\,f[X-(1-u)/z]$ with modes $\Hop^u_m$
\cite[Def.~1.1]{Q99}.  For a nonzero $F\in\Q[t^{\pm}]$ (or $F\in\Lam\otimes\Q[t^{\pm}]$)
we write $\ord F$ for the order of vanishing at $t=-1$; for a symmetric function we take
the minimum over the coefficients in a fixed basis, which is basis-independent.

\paragraph{(A), the ribbon-side $(1+t)$.}  On $ts=1$ the two-bead sector of
$[R_e(t),R_f(s)]$ vanishes identically \cite[Cor.~4.2(iii)]{Q96}, so the commutator is
carried entirely by its one-bead sector, whose matrix elements are divisible by $(1+t)$.

\paragraph{(B), the vertex-side $(1+t)$.}  For $s=t^{-1}$,
\begin{equation}\label{eq:B}
  \Hop^{t}_m \Hop^{1/t}_n-t^{-1}\Hop^{1/t}_n \Hop^{t}_m
  =\bigl(1-t^{-1}\bigr)t^{-m}\,h_{m+n}\bigl[(1+t)X\bigr]\cdot,
\end{equation}
a multiplication operator, the plethysm being $p_k\mapsto(1+t^k)p_k$ \cite[Thm.~3.4]{Q99}.

\paragraph{The trap.}  Both objects vanish at $t=-1$ in \emph{some} instances, and
$(1+t^k)$ is the Hall--Littlewood/Schur-$Q$ doubling deformation, so a common origin is
plausible.  That is a shape match.  Agreement of values at $t=-1$ is worthless here: it is
a check with a kernel in exactly the direction it is meant to test.  We therefore fix in
advance the two coordinates that decide the question, neither of which can be tuned:
\begin{enumerate}
\item[(S1)] the \emph{order} of vanishing at $t=-1$, as a function on each side's own
parameter space (not the value, and not the fact of vanishing);
\item[(S2)] the origin of the alphabet: is $(1+t)$ a contraction (Wick) constant of a
$t$-twisted fermion, or something else?
\end{enumerate}
A merge and a split are equally good outcomes.  ``Both have a $(1+t)$'' is not an outcome.

\section{The ribbon side: one integer controls everything}

We start from the two-parameter matrix elements, which are already proved:

\begin{proposition}[{\cite[Thm.~4.1(1)]{Q96}}]\label{prop:q96}
Let $e,f\ge1$, let $M$ be a Maya set and let $M'=M\cdot(a,e+f)$ be a legal $(e{+}f)$-move.
Put
\[
 N=\#\bigl(M\cap(a,a+e+f)\bigr),\quad A=\#\bigl(M\cap(a,a+f)\bigr),\quad
 B=\#\bigl(M\cap(a,a+e)\bigr),
\]
$m_e=\mathbf 1[a+e\in M]$ and $m_f=\mathbf 1[a+f\in M]$.  Then
\[
  \langle M'|[R_e(t),R_f(s)]|M\rangle
  = s^{A}t^{\,N-1-A}\bigl(t-m_f(1+t)\bigr)-t^{B}s^{\,N-1-B}\bigl(s-m_e(1+s)\bigr).
\]
\end{proposition}

\begin{definition}\label{def:kappa}
With the notation of Proposition~\ref{prop:q96} set
\[
  \boxed{\ \kappa=\kappa^{e,f}_{a}(M):=2\bigl(A+B-N\bigr)+m_e+m_f\ },
  \qquad P:=N-2A-m_f .
\]
\end{definition}

$\kappa$ is an integer attached to the bead configuration; it is not a fitted parameter and
it is not constant --- on the sweep of \S\ref{sec:verif} it takes every value in
$\{-4,\dots,6\}$.

\begin{theorem}[the hyperbola closed form]\label{thm:A}
Let $e,f\ge1$ and $s=t^{-1}$.  Then for every Maya set $M$ and every legal
$(e{+}f)$-move $(a,e+f)$, with $M'=M\cdot(a,e+f)$,
\[
  \boxed{\ \langle M'|[R_e(t),R_f(t^{-1})]|M\rangle
  \;=\;(-1)^{N}\,(-t)^{P}\,\bigl(1-(-t)^{\kappa}\bigr).\ }
\]
Consequently:
\begin{enumerate}
\item[(i)] the matrix element is zero if and only if $\kappa=0$;
\item[(ii)] if $\kappa\neq0$ its order of vanishing at $t=-1$ is \textbf{exactly $1$}, with
\[
  \langle M'|[R_e(t),R_f(t^{-1})]|M\rangle\;\equiv\;(-1)^{N}\kappa\,(t+1)
  \pmod{(t+1)^{2}} ;
\]
\item[(iii)] in particular $(1+t)$ divides every one-bead matrix element, uniformly in
$e,f$ and in $M$, and divides no such element twice.
\end{enumerate}
\end{theorem}

\begin{proof}
Since $m_f\in\{0,1\}$ we have $t-m_f(1+t)=t$ when $m_f=0$ and $=-1$ when $m_f=1$, i.e.
\[
  t-m_f(1+t)=(-1)^{m_f}\,t^{\,1-m_f}.
\]
Likewise $s-m_e(1+s)=(-1)^{m_e}s^{\,1-m_e}$, which at $s=t^{-1}$ is
$(-1)^{m_e}t^{\,m_e-1}$.  Substituting $s=t^{-1}$ into Proposition~\ref{prop:q96},
\[
  s^At^{N-1-A}=t^{\,N-1-2A},\qquad t^Bs^{N-1-B}=t^{\,2B-N+1},
\]
so the matrix element equals
\begin{equation}\label{eq:binom}
  \Xi:=(-1)^{m_f}\,t^{\,N-2A-m_f}\;-\;(-1)^{m_e}\,t^{\,2B-N+m_e}
     \;=\;(-1)^{m_f}t^{P}-(-1)^{m_e}t^{Q},
\end{equation}
where $P=N-2A-m_f$ and $Q:=2B-N+m_e$.  Note
\[
  Q-P=2B-N+m_e-N+2A+m_f=2(A+B-N)+m_e+m_f=\kappa .
\]
Now the two signs collapse.  Since $P\equiv N+m_f$ and $Q\equiv N+m_e \pmod 2$,
\[
  (-1)^{m_f}=(-1)^{N+P},\qquad (-1)^{m_e}=(-1)^{N+Q},
\]
so \eqref{eq:binom} becomes
\[
  \Xi=(-1)^{N}\bigl[(-1)^{P}t^{P}-(-1)^{Q}t^{Q}\bigr]
     =(-1)^{N}\bigl[(-t)^{P}-(-t)^{Q}\bigr]
     =(-1)^{N}(-t)^{P}\bigl(1-(-t)^{\kappa}\bigr),
\]
using $Q=P+\kappa$.  This is the closed form.

For (i): $(-t)^P$ is a unit in $\Q[t^{\pm}]$, so $\Xi=0$ iff $(-t)^{\kappa}=1$ in
$\Q[t^{\pm}]$, iff $\kappa=0$.  For (ii): substituting $t=-1$ gives $(-t)=1$ and
$1-1^{\kappa}=0$, so $(t+1)\mid\Xi$ always.  Differentiating,
$\frac{d}{dt}\bigl(1-(-t)^{\kappa}\bigr)=\kappa(-t)^{\kappa-1}$, which at $t=-1$ equals
$\kappa$; and $(-1)^N(-t)^P$ equals $(-1)^N$ at $t=-1$.  Hence
$\Xi'(-1)=(-1)^{N}\kappa$, nonzero precisely when $\kappa\neq0$, which is (ii); and (iii)
follows.
\end{proof}

\begin{remark}[what the $(1+t)$ \emph{is}, on this side]\label{rem:cyclotomic}
Theorem~\ref{thm:A} exhibits the ribbon-side $(1+t)$ as the first cyclotomic factor of a
quantum integer at $q=-t$: for $\kappa>0$,
\[
  1-(-t)^{\kappa}=\bigl(1-(-t)\bigr)\sum_{j=0}^{\kappa-1}(-t)^{j}=(1+t)\,[\kappa]_{-t}.
\]
It is a \emph{divisor}, and the exponent $\kappa$ is a bead statistic.  This is the same
shape as the two-bead sector, whose matrix elements carry the factor $1-(ts)^{k}$
\cite[Thm.~4.1(2)]{Q96} and which is what makes that sector vanish on $ts=1$; the one-bead
sector is the same mechanism with $ts$ replaced by $-t$ and the interference index $k$
replaced by $\kappa$.  Nothing here is an alphabet, and no companion factor $(1+t^{j})$
with $j\ge2$ ever appears.
\end{remark}

\begin{corollary}\label{cor:Aorder}
On the hyperbola $ts=1$ the operator $[R_e(t),R_f(t^{-1})]$ has bead-number exactly $1$,
and the order function
\[
  \omega_A:\ (e,f,M,M')\ \longmapsto\ \ord\ \langle M'|[R_e(t),R_f(t^{-1})]|M\rangle
  \ \in\ \{1\}
\]
is \textbf{constant} on the whole of its (nonvanishing) domain.
\end{corollary}

\begin{proof}
The two-bead sector vanishes identically on $ts=1$ because each of its summands carries the
factor $1-(ts)^{k}$ \cite[Thm.~4.1(2), Cor.~4.2(iii)]{Q96}, and sectors of bead-number
$\ge3$ do not occur since each $R_g$ has bead-number $1$ \cite[Conv.~2.1]{Q92}.  The
constancy is Theorem~\ref{thm:A}(ii).
\end{proof}

\section{The vertex side: palindromy fixes the parity of the order}

\begin{lemma}[the $h$-expansion]\label{lem:hexp}
For $N\ge0$, \quad $h_N\bigl[(1+t)X\bigr]=\sum_{k=0}^{N}t^{\,N-k}\,h_k\,h_{N-k}$.
\end{lemma}

\begin{proof}
$\sigma[(1+t)Xz]=\sigma[Xz]\,\sigma[tXz]$ because $p_n[(1+t)X]=p_n[X]+p_n[tX]$ and
$\sigma[A+B]=\sigma[A]\sigma[B]$.  Comparing coefficients of $z^N$ and using
$h_j[tX]=t^{j}h_j$ gives the claim.
\end{proof}

\begin{lemma}[palindromy]\label{lem:palin}
Fix partitions $\lambda,\nu$ with $|\nu|=|\lambda|+N$, and set
$c_k:=\bigl\langle s_\nu,\ h_kh_{N-k}\,s_\lambda\bigr\rangle\in\Z_{\ge0}$ and
$\Pi_{\nu\lambda}(t):=\bigl\langle s_\nu,\ h_N[(1+t)X]\,s_\lambda\bigr\rangle$.
Then $\Pi_{\nu\lambda}(t)=\sum_{k=0}^{N}c_kt^{\,N-k}$ and $c_k=c_{N-k}$, i.e.\
$\Pi_{\nu\lambda}(t)=t^{N}\Pi_{\nu\lambda}(1/t)$.
\end{lemma}

\begin{proof}
The displayed expansion is Lemma~\ref{lem:hexp} paired against $s_\nu$.  The symmetry is
$h_kh_{N-k}=h_{N-k}h_k$, which gives $c_k=c_{N-k}$ on the nose.
\end{proof}

\begin{lemma}[order parity of a self-reciprocal polynomial]\label{lem:selfrec}
Let $0\neq \Pi\in\Q[t]$ satisfy $\Pi(t)=t^{N}\Pi(1/t)$.  Then the multiplicity of $t=-1$
as a root of $\Pi$ is congruent to $N$ modulo $2$.
\end{lemma}

\begin{proof}
Write $\Pi=t^{\alpha}Q$ with $Q(0)\neq0$.  From $\Pi(t)=t^N\Pi(1/t)$ we get
$Q(t)=t^{\,N-2\alpha}Q(1/t)$, and since $Q(0)\neq0$ this forces $\deg Q=N-2\alpha$.
Factor $Q=(t-1)^{m_+}(t+1)^{m_-}R$ with $R(\pm1)\neq0$.  The nonreal-or-non-unit roots of
$Q$ occur in pairs $\{\rho,1/\rho\}$ with equal multiplicities and $\rho\neq\pm1$, so
$\deg R$ is even and $\deg Q\equiv m_++m_-\pmod2$.  Finally $m_+$ is even: writing
$Q=(t-1)^{m_+}S$ with $S(1)\neq0$, self-reciprocity of $Q$ forces $S=(-1)^{m_+}S^{*}$ where
$S^{*}(t)=t^{\deg S}S(1/t)$; if $m_+$ were odd this would give $S(1)=-S^{*}(1)=-S(1)$,
hence $S(1)=0$, a contradiction.  So
$m_-\equiv\deg Q=N-2\alpha\equiv N\pmod2$.
\end{proof}

\begin{theorem}[the vertex-side order function]\label{thm:B}
Let $N=m+n\ge0$ and let $D_{m,n}(t)=(1-t^{-1})t^{-m}h_{N}[(1+t)X]$ be the defect operator
of \eqref{eq:B}.  Then:
\begin{enumerate}
\item[(i)] $\ord\,D_{m,n}=\ord\,h_N[(1+t)X]=N\bmod 2$, as an element of
$\Lam\otimes\Q[t^{\pm}]$;
\item[(ii)] every nonzero Schur matrix element of $D_{m,n}$ has order of vanishing
$\equiv N\pmod 2$; in particular for $N$ \emph{even} no matrix element has order $1$;
\item[(iii)] for $N$ even the order $0$ is attained: with $\lambda=\varnothing$ and
$\nu=(N)$ one has $\Pi_{\nu\lambda}(t)=1+t+\cdots+t^{N}$, whose value at $t=-1$ is $1$.
\end{enumerate}
Hence the order function
$\omega_B:(\lambda,\nu,m,n)\mapsto\ord\langle s_\nu|D_{m,n}|s_\lambda\rangle$ takes both
parities and attains the value $0$.
\end{theorem}

\begin{proof}
(i) The prefactor $(1-t^{-1})$ takes the value $2$ at $t=-1$ and $t^{-m}$ is a unit, so
neither changes the order.  In the power-sum basis
$h_N[(1+t)X]=\sum_{\mu\vdash N}z_\mu^{-1}\prod_i(1+t^{\mu_i})p_\mu$, and $(1+t^{j})$
vanishes at $t=-1$ exactly when $j$ is odd, to order $1$.  So the coefficient of $p_\mu$ has
order equal to the number $o(\mu)$ of odd parts of $\mu$, and the order of the whole is
$\min_{\mu\vdash N}o(\mu)$.  Since $o(\mu)\equiv|\mu|=N\pmod2$, this minimum is $0$ for $N$
even (take $\mu=(2^{N/2})$) and $1$ for $N$ odd (take $\mu=(2^{(N-1)/2},1)$).

(ii) Combine Lemmas~\ref{lem:palin} and~\ref{lem:selfrec}.

(iii) $\langle s_{(N)},h_kh_{N-k}\rangle=K_{(N),(k,N-k)}=1$ for every $k$, so $c_k=1$ for
all $k$ and $\Pi=\sum_{k=0}^Nt^{k}$; at $t=-1$ this alternating sum is $1$ when $N$ is even.
\end{proof}

\begin{remark}[a ladder, not a single slice]\label{rem:ladder}
The proof of (i) shows more: the per-partition orders form the full ladder
$\{N\bmod2,\ N\bmod 2+2,\ \dots,\ N\}$, and the Schur matrix elements realise the same
ladder (computed in \S\ref{sec:verif}).  So on the vertex side the order at $t=-1$ is a
genuine \emph{grading} by the number of odd parts.  On the ribbon side there is no ladder:
the order is $1$ or the element is zero.
\end{remark}

\section{The decision}

\begin{theorem}[Q105]\label{thm:decision}
The objects (A) and (B) are \textbf{distinct}.  Precisely: there is no map $\phi$ from the
parameter set of (A) to that of (B) with $\omega_A=\omega_B\circ\phi$ and $\phi$ surjective
onto the parameter $m+n$ of (B).
\end{theorem}

\begin{proof}
By Corollary~\ref{cor:Aorder}, $\omega_A\equiv1$.  By Theorem~\ref{thm:B}(ii),
$\omega_B\equiv m+n\pmod 2$.  If $\omega_A=\omega_B\circ\phi$ then $m+n$ is odd at every
point of the image of $\phi$; but $D_{m,n}\neq0$ for every $m+n\ge0$ (its $\lambda=\varnothing$,
$\nu=(m+n)$ matrix element is $(1-t^{-1})t^{-m}[m+n+1]_t\neq0$), so the even-$(m+n)$ half of
(B) is a nonempty family of nonzero operators outside the image.  Hence $\phi$ is not
surjective.
\end{proof}

\begin{remark}[what the shape match confused]\label{rem:divisor-vs-alphabet}
The two $(1+t)$'s play different grammatical roles.
\begin{center}
\begin{tabular}{lll}
 & \textbf{(A) ribbon} & \textbf{(B) vertex}\\\hline
role of $(1+t)$ & a \emph{divisor} of the object & a \emph{plethystic alphabet} $X\mapsto(1+t)X$\\
divides the object? & always & only when $m+n$ is odd\\
order at $t=-1$ & exactly $1$, always & $\equiv m+n \pmod 2$, ladder up to $m+n$\\
higher modes $(1+t^{j})$, $j\ge2$ & never occur & all occur\\
bead-number of the operator & exactly $1$ & unbounded (multiplication)
\end{tabular}
\end{center}
``Divisible by $(1+t)$'' and ``obtained by the plethysm $(1+t)X$'' are different predicates
that agree only in odd degree.  The tempting inference identified a factor with an argument.
\end{remark}

\begin{remark}[the separator would have fired the other way too]
Nothing above was tuned to produce a split.  Had $\omega_A$ come out equal to
$e+f \bmod 2$, the dictionary $m+n=e+f$ would have matched the two order functions exactly
and the honest report would have been a merge.  It came out constant.
\end{remark}

\section{The alphabet: (1-t) is the twist, (1+t) is the pole}\label{sec:alphabet}

The remaining question is (S2): where does the $(1+t)$ of \eqref{eq:B} come from?  The
published twist alphabet for the Hall--Littlewood boson--fermion correspondence is
$(1-t)$, not $(1+t)$: Necoechea--Rozhkovskaya obtain the Hall--Littlewood presentation from
the Schur one ``by a simple twist of the fields of charged free fermions by multiplication
by $E(-u/t)$ or $H(u/t)$'', with $\Psi^{+}(u)=E(-u/t)\Phi^{+}(u)$ and
$\Psi^{-}(u)=H(u/t)\Phi^{-}(u)$; the twisted fields satisfy an anticommutation relation
whose right-hand side is $\delta(u,v)(1-t)^{2}$, and
$H(u)E(-u/t)=\exp\bigl(\sum_{n\ge1}\tfrac{(1-t^{n})p_n}{n}u^{-n}\bigr)$
\cite[ll.~385, 392--394, 398--405, 735, 764]{NR}.  Note $t\mapsto-t$ does \emph{not} carry
$1-t^{n}$ to $1+t^{n}$: it does so only for odd $n$.  So either contracting a
$(1-t)$-twisted field against a $(1-t^{-1})$-twisted one produces $(1+t)$ after
simplification, or the two alphabets are genuinely different.  We show the latter, and
identify what $(1+t)$ actually is.

\begin{definition}\label{def:htilde}
For parameters $a,b$ put
$\ \Htil_{a,b}(z)f[X]=\sigma[Xz]\ f\!\left[X-\dfrac{a-b}{z}\right]$,
the shift being the plethystic substitution $p_n\mapsto p_n-(a^n-b^n)z^{-n}$.
The Hall--Littlewood operator is the pin $a=1$: $\Hop_u=\Htil_{1,u}$.
\end{definition}

\begin{proposition}[zero and pole]\label{prop:zeropole}
Let $\Phi'$ denote the normal-ordered series
$\Phi'=\sigma[Xz_1]\sigma[Xz_2]\ f\bigl[X-\tfrac{a-b}{z_1}-\tfrac{c-d}{z_2}\bigr]$.  Then
\[
  (z_1-b z_2)\,\Htil_{a,b}(z_1)\Htil_{c,d}(z_2)=(z_1-a z_2)\,\Phi',
  \qquad
  (z_2-d z_1)\,\Htil_{c,d}(z_2)\Htil_{a,b}(z_1)=(z_2-c z_1)\,\Phi' .
\]
In particular the contraction constant of $\Htil_{a,b}(z_1)$ against $\Htil_{c,d}(z_2)$ is
$\dfrac{z_1-az_2}{z_1-bz_2}$: it has its \emph{zero} at $z_1=az_2$ and its \emph{pole} at
$z_1=bz_2$, and its alphabet is $(a-b)$.
\end{proposition}

\begin{proof}
The only contraction is between the shift of the left operator and the $\sigma[Xz_2]$ of the
right one.  Substituting $X\mapsto X-\frac{a-b}{z_1}$ into $\sigma[Xz_2]$ gives
$\sigma[Xz_2]\cdot\sigma\bigl[-(a-b)\tfrac{z_2}{z_1}\bigr]$, and
\[
  \sigma\Bigl[-(a-b)\tfrac{z_2}{z_1}\Bigr]
  =\exp\Bigl(-\sum_{n\ge1}\tfrac{a^n-b^n}{n}\bigl(\tfrac{z_2}{z_1}\bigr)^{n}\Bigr)
  =\frac{1-a z_2/z_1}{1-b z_2/z_1}=\frac{z_1-az_2}{z_1-bz_2},
\]
the middle step being $-\sum_n\frac{(a^n-b^n)y^n}{n}=\log(1-ay)-\log(1-by)$.  Moving the
shift of the left operator to the right of $\sigma[Xz_2]$ therefore produces exactly this
factor times $\Phi'$, which is the first identity; the second is the same computation with
the roles of the two operators exchanged.  (At $a=c=1$, $b=t$, $d=s$ the two identities give
$(z_1-tz_2)\Hop_t(z_1)\Hop_s(z_2)=(sz_1-z_2)\Hop_s(z_2)\Hop_t(z_1)$, which is
\cite[Thm.~2.4]{Q99}.)
\end{proof}

\begin{theorem}[the defect alphabet is the pole, not the twist]\label{thm:C}
Fix $b\neq0$ and let $c=1$, $d=a/b$, with $a$ \emph{free}.  Write $\Phi'_{l,l'}$ for the
coefficient of $z_1^{l}z_2^{l'}$ in $\Phi'$.  Then for every $f\in\Lam$ and every $N$,
\[
  \boxed{\ \sum_{l\in\Z}b^{\,l}\,\Phi'_{l,\,N-l}\,f\;=\;h_N\bigl[(1+b)X\bigr]\cdot f\ }
\]
(the sum being finite, and the right side $0$ for $N<0$).  The right-hand side does not
involve $a$, whereas the two twist alphabets are $(a-b)$ and $(1-a/b)$ and both do.
\end{theorem}

\begin{proof}
Two things happen on the diagonal $z_1=bz_2$, which is the pole of the contraction of
Proposition~\ref{prop:zeropole}.

\emph{The shift cancels.}  The total plethystic shift in $\Phi'$ is
$-\frac{a-b}{z_1}-\frac{c-d}{z_2}$, i.e.\ $p_n\mapsto p_n-\frac{a^n-b^n}{z_1^{n}}
-\frac{c^n-d^n}{z_2^{n}}$.  At $z_1=bz_2$, $c=1$ and $d=a/b$ the bracket is
\[
  \frac{1}{z_2^{n}}\Bigl[-\frac{a^{n}-b^{n}}{b^{n}}-\bigl(1-(a/b)^{n}\bigr)\Bigr]
  =\frac{1}{z_2^{n}}\Bigl[-\frac{a^{n}}{b^{n}}+1-1+\frac{a^{n}}{b^{n}}\Bigr]=0
\]
for every $n\ge1$ and every $a$.  So $f\bigl[X-\frac{a-b}{z_1}-\frac{c-d}{z_2}\bigr]$
restricted to the diagonal is $f[X]$, and $f$ passes through untouched.

\emph{The creations double.}  $\sigma[Xz_1]\sigma[Xz_2]$ at $z_1=bz_2$ is
$\sigma[X\,bz_2]\sigma[Xz_2]=\sigma\bigl[(1+b)Xz_2\bigr]$, since
$\sigma[A]\sigma[B]=\sigma[A+B]$ and $p_n[Xbz_2]+p_n[Xz_2]=(1+b^{n})p_nz_2^{n}$.

Now the diagonal sum $\sum_l b^{l}\Phi'_{l,N-l}$ is precisely coefficient extraction of
$z_2^{N}$ from $\Phi'|_{z_1=bz_2}$: writing $\Phi'=\sum_{l,l'}\Phi'_{l,l'}z_1^{l}z_2^{l'}$
and setting $z_1=bz_2$ gives $\sum_{l,l'}\Phi'_{l,l'}b^{l}z_2^{l+l'}$.  The substitution is
legitimate because for fixed $f$ of degree $\deg f=D$ the shift is a Laurent polynomial with
$\Phi'_{l,l'}f=0$ unless $l\ge-D$ and $l'\ge-D$, so each coefficient of $z_2^{N}$ is a
finite sum (this is the finiteness argument of \cite[Lem.~3.3]{Q99}, unchanged).  Combining
the two displays, that coefficient is $h_N[(1+b)X]\,f$.
\end{proof}

\begin{corollary}[answer to (S2): option (ii)]\label{cor:S2}
$(1+t)$ is not a contraction constant of the $t$-twisted fields.  The contraction constant
is $\frac{z_1-z_2}{z_1-tz_2}$, whose alphabet is $(1-t)$ --- exactly the published
Necoechea--Rozhkovskaya twist --- and which has a \emph{pole}, not a value, at the point
$z_1=tz_2$ where the $(1+t)$ appears.  The $(1+t)$ is $1+{}$(pole location): it is the
doubled creation alphabet $\sigma[Xz_1]\sigma[Xz_2]$ restricted to the diagonal.  The two
alphabets are functions of different data --- the twist needs the zero \emph{and} the pole,
the defect needs only the pole --- and they coincide in form only because the
Hall--Littlewood normalisation pins the zero at $a=1$.
\end{corollary}

\begin{proof}
Theorem~\ref{thm:C}: hold the pole $b=t$ fixed and vary the zero $a$.  The twist alphabets
$(a-b)$ and $(1-a/b)$ move; the defect alphabet $(1+b)$ does not.  A quantity that is
constant along a deformation cannot equal one that varies along it.
\end{proof}

\begin{remark}[a second, cheaper witness]\label{rem:tminus}
Inside the pinned family $a=1$, where both alphabets are functions of the single parameter
$t$, the same conclusion follows from $t\mapsto-t$.  Equation \eqref{eq:B} holds for every
value of the parameter, so applying it at $-t$ gives
\[
  \Hop^{-t}_m\Hop^{-1/t}_n+t^{-1}\Hop^{-1/t}_n\Hop^{-t}_m
  =\bigl(1+t^{-1}\bigr)(-t)^{-m}\,h_{m+n}\bigl[(1-t)X\bigr]\cdot,
\]
verified $252/252$ in \S\ref{sec:verif}.  The defect alphabet is now $(1-t)$ and the twist
alphabet is $(1+t)$: in every odd mode the two have exchanged.  A pair of quantities that
swap under a reparametrisation of the family is not one quantity.
\end{remark}

\begin{remark}[on the standing hypothesis I was asked not to build on]
Rick's Day~180 letter asserts at $75\%$ stated confidence that $(1+t)X$ ``is exactly the
Wick constant of a $t$-twisted charged fermion against its dual''.  Corollary~\ref{cor:S2}
refutes the assertion as stated, and Proposition~\ref{prop:zeropole} says exactly how: he
identified the correct \emph{locus} --- the $(1+t)$ does live at the contraction's
singularity --- but not the correct \emph{object}, since at that locus the contraction is
infinite.  The relationship is ``$1+{}$pole of the Wick constant'', not ``value of the Wick
constant''.  This is the same failure mode as
\texttt{a-confirmed-operator-is-not-a-confirmed-form}: right which, wrong how.
\end{remark}

\section{Verification}\label{sec:verif}

All primary computations for this paper were run on two engines written this session:
\texttt{clio-vega/proofs@d9044b7:2026-09-08-c2-Q105-code/abacus.py} (Maya sets and bead
moves only; no Young diagrams, no symmetric functions) and \texttt{.../vertex.py} (power-sum
dictionaries; no $h$-basis, no abacus); that directory's \texttt{README.md} tabulates which
script establishes which line below.  The two engines of the 8 September self-review,
\texttt{clio-vega/rick-review@314f24a:2026-09-08-selfreview-code/ribbon.py} (border strips on
Young diagrams) and \texttt{hbasis.py} ($h$-basis dictionaries, no power sums), were used only as
\emph{cross-checks} and share no code path with either primary.

\begin{enumerate}
\item \textbf{Primary engine calibrated against a proved theorem.}  All one-bead matrix
elements of $[R_e(t),R_f(s)]$ at generic symbolic $(t,s)$, for $|\lambda|\le6$ and
$1\le e,f\le4$: $1232/1232$ agree with Proposition~\ref{prop:q96}; $597$ two-bead entries
seen.  Untuned anchor: $[R_1,R_2]s_\varnothing=(1+t)s_{(2,1)}$, recomputed from the abacus
and matching \cite[\S3]{Q92}.

\item \textbf{Theorem~\ref{thm:A}, closed form.}  $2588/2588$ over the same range at
$s=1/t$ --- \emph{every} legal $(e{+}f)$-move, not only the nonzero ones.  The split is
$1232$ nonzero and $1356$ zero, and the predicate $\kappa=0$ separates them with $0$
failures in either direction.  $\kappa$ realises every value in $\{-4,\dots,6\}$, so the
statement is not vacuous in its own parameter.

\item \textbf{Theorem~\ref{thm:A}(ii), order.}  All $1232$ nonzero elements have order of
vanishing exactly $1$ at $t=-1$; the histogram is $\{1:1232\}$ with no other order
occurring.  The leading coefficient test $\Xi'(-1)=(-1)^{N}\kappa$ passes $1232/1232$.

\item \textbf{Two-bead sector on the hyperbola.}  $0$ of the $597$ generically-nonzero
two-bead entries survive at $s=1/t$, confirming \cite[Cor.~4.2(iii)]{Q96} independently.
(We do not use \cite[Cor.~4.2(iv)]{Q96}, which carries a known correction for $e\le2$
recorded in the annotation to that paper.)

\item \textbf{Cross-check of Theorem~\ref{thm:A} on \texttt{ribbon.py}.}  $764/764$ one-bead
elements of order exactly $1$, with the same leading coefficient $(-1)^{N}\kappa$ and $0$
failures.  That engine computes $R_g$ by enumerating border strips on Young diagrams and
knows nothing of $\kappa$.

\item \textbf{Lemma~\ref{lem:hexp}.}  The power-sum route
($h_N=\sum_{\mu}p_\mu/z_\mu$ then $p_k\mapsto(1+t^k)p_k$) and the creation route
($\sum_k t^{N-k}h_kh_{N-k}$) agree coefficientwise for $N=0,\dots,8$.

\item \textbf{Theorem~\ref{thm:B}.}  $\ord h_N[(1+t)X]=N\bmod2$ for $N=0,\dots,10$, with the
per-partition ladder $\{N\bmod2,\dots,N\}$ realised in full each time.  All $513$ nonzero
Schur matrix elements with $N\le6$ and $|\lambda|\le4$ are palindromic ($0$ exceptions), and
their order histogram is $\{0,2,4,\dots\}$ for $N$ even and $\{1,3,5,\dots\}$ for $N$ odd:
\emph{no matrix element of even $N$ has order $1$}.  Explicit order-$0$ witnesses at $N=2$:
$\langle s_{(2)}|\cdot|s_\varnothing\rangle=1+t+t^{2}$ and
$\langle s_{(1,1)}|\cdot|s_\varnothing\rangle=t$.

\item \textbf{Cross-check of Theorem~\ref{thm:B} on \texttt{hbasis.py}.}  Recomputing the
defect \eqref{eq:B} from the $h$-basis engine for $0\le m,n\le4$ gives coefficient orders
$0,1,0,1,0,1,0,1,0$ for $m+n=0,\dots,8$: exactly $N\bmod2$.

\item \textbf{Equation \eqref{eq:B} re-derived.}  Theorem~\cite[3.4]{Q99} was recomputed
from scratch on \texttt{vertex.py}: $252/252$ at $u=t$ ($|\mu|\le3$, $-2\le m,n\le3$), of
which $182$ have nonzero defect.  The same identity at $u=-t$: $252/252$, giving
Remark~\ref{rem:tminus}.

\item \textbf{Proposition~\ref{prop:zeropole}} with \emph{four} independent symbols
$a,b,c,d$: $64+64$ agreements, $0$ mismatches.

\item \textbf{Theorem~\ref{thm:C}}, the varying test: with $b$ symbolic and $a$ a free
symbol, the diagonal sum equals $h_N[(1+b)X]$ in $35/35$ cases ($|\mu|\le3$,
$-1\le N\le3$).  Specialising $a\in\{1,-1,2,\tfrac12\}$ at fixed pole $b=t$ moves the twist
alphabets to $(1-t,\,1-t^2,\,1-t^3)$, $(-1-t,\,1-t^2,\,-1-t^3)$, $(2-t,\,4-t^2,\,8-t^3)$,
$(\tfrac12-t,\,\tfrac14-t^2,\,\tfrac18-t^3)$ while the defect alphabet stays
$(1+t,\,1+t^2,\,1+t^3)$ throughout.
\end{enumerate}

\section{What is not claimed}

\begin{enumerate}
\item \textbf{Not claimed: that the two objects are unrelated.}  They are related, and
Theorem~\ref{thm:C} says how: both are governed by the same distinguished point.  On the
vertex side that point is the pole $z_1=tz_2$ of the contraction; on the ribbon side the
$-t$ of $1-(-t)^{\kappa}$ plays the analogous role for the height statistic.  What is
refuted is that they are the \emph{same} $(1+t)$, i.e.\ that one derivation transports to
the other.

\item \textbf{Not claimed: a dictionary between $(e,f)$ and $(k,m+n)$.}  The brief asked
whether any dictionary makes the order functions agree.  The answer is that a dictionary can
only do so by landing entirely in odd $m+n$ (Theorem~\ref{thm:decision}), and no principle
selects the odd half of a family that is defined for all degrees.  If a dictionary is ever
produced it must explain what pins that parity; I have none.

\item \textbf{Not claimed: anything about the two-bead sector off $ts=1$.}  Everything here
is on the hyperbola.

\item \textbf{Not claimed: Rick's \S2 residue suggestion.}  Step~3 of the brief (is
Theorem~\eqref{eq:B} literally a formal residue at $u=tz_2/z_1=1$?) was not attempted.
Proposition~\ref{prop:zeropole} makes it look very likely --- the pole is simple and the
$\delta$ is its residue --- but ``looks like a residue'' is precisely the kind of statement
this paper exists to refuse.

\item \textbf{On the Necoechea--Rozhkovskaya citation.}  I quote the twist alphabet
$(1-t^{n})$ and the one-sided, two-dressing form of the twist.  I do \emph{not} claim that
paper says anything about $(1+t)$, about the hyperbola $st=1$, or about my defect; it does
not, and Corollary~\ref{cor:S2} is mine.
\end{enumerate}

\begin{thebibliography}{9}
\bibitem{Q92} Clio, \emph{The cross-rank ribbon commutator}, 7 September 2026,
\texttt{proofs/2026-09-07-Q92-cross-rank-commutator.tex}.
\bibitem{Q96} Clio, \emph{Infinite order over the sublattice subring}, 7 September 2026
(cycle 2), \texttt{proofs/2026-09-07-c2-Q96-order-over-sublattice.tex}; Theorem~4.1 and
Corollary~4.2, with the annotation of 8 September 2026 to Corollary~4.2(iv).
\bibitem{Q99} Clio, \emph{The two-parameter Hall--Littlewood exchange relation, and what
actually lives on the hyperbola $st=1$}, 8 September 2026,
\texttt{proofs/2026-09-08-Q99-two-parameter-exchange.tex}.
\bibitem{NR} A.~Necoechea and N.~Rozhkovskaya, \emph{Boson-fermion correspondence for
Hall-Littlewood polynomials revisited}, arXiv:1902.10049; LaTeX source read at source
8 September 2026, lines 385, 392--394, 398--405, 735, 764.
\end{thebibliography}

\end{document}
